\documentclass{article} % For LaTeX2e
\usepackage{iclr2026_conference,times}


% Minimal math_commands stub so that % Minimal math_commands stub so that % Minimal math_commands stub so that \input{math_commands.tex} resolves.
% The manuscript only relies on standard amsmath / amssymb macros.
\newcommand{\eps}{\varepsilon}
 resolves.
% The manuscript only relies on standard amsmath / amssymb macros.
\newcommand{\eps}{\varepsilon}
 resolves.
% The manuscript only relies on standard amsmath / amssymb macros.
\newcommand{\eps}{\varepsilon}


\usepackage{hyperref}
\usepackage{url}
\usepackage[utf8]{inputenc}
\usepackage[T1]{fontenc}
\usepackage{microtype} % 改善字间距


\usepackage{amsmath, amsfonts, amssymb, amsthm}
\usepackage{nicefrac} % 漂亮的斜分式

\newtheorem{theorem}{Theorem}[section]
\newtheorem{proposition}[theorem]{Proposition}
\newtheorem{lemma}[theorem]{Lemma}
\newtheorem{definition}[theorem]{Definition}
\newtheorem{remark}[theorem]{Remark}


\usepackage{algorithm}
\usepackage{algpseudocode}


\usepackage[table]{xcolor} % 必须在这里带上 [table] 选项，防止后续冲突
\usepackage{graphicx}
\usepackage{float}
\usepackage{placeins}

\definecolor{headerpink}{RGB}{241,222,218}
\definecolor{rowgray}{RGB}{244,244,244}
\definecolor{mydarkblue}{rgb}{0,0.4,0.8} 
\definecolor{mylightblue}{rgb}{0.5,0.75,1}
\definecolor{NavyBlue}{HTML}{000080}


% ==============================================================================
% 5. 表格工具
% ==============================================================================
\usepackage{booktabs}
\setlength{\heavyrulewidth}{\lightrulewidth} % 使表格顶线/底线与中线等宽 (根据你之前的需求)
\setlength{\aboverulesep}{0pt}               % 消除彩色行白边
\setlength{\belowrulesep}{0pt}


\usepackage{enumitem} % 列表缩进控制
\usepackage{pifont}   % \ding 符号
\usepackage{xspace}   % 智能空格
\usepackage{url}

% 批注命令
\newcommand{\XL}[1]{\textcolor{orange}{[XL: #1]}}

% 方法名 CHAIN 宏定义
\DeclareRobustCommand{\method}{%
  \texorpdfstring{\ifmmode\text{\normalfont\ttfamily\bfseries CHAIN}\else\texttt{\textbf{CHAIN}}\xspace\fi}{CHAIN}%
}

\usepackage{hyperref}
\hypersetup{
    colorlinks=true,
    linkcolor=red,
    citecolor=mydarkblue,
    filecolor=magenta,
    urlcolor=magenta,
}


\title{Perturbation Reliability Score: Training-Free \\ Uncertainty Quantification for LLM Reasoning}

% Anonymous for double-blind review (ICLR 2026).
\author{Anonymous authors\\
Paper under double-blind review}

\newcommand{\fix}{\marginpar{FIX}}
\newcommand{\new}{\marginpar{NEW}}

%\iclrfinalcopy % Uncomment for camera-ready version, but NOT for submission.
\begin{document}


\maketitle

\begin{abstract}
Large language models can produce fluent reasoning traces that are nonetheless
wrong, and reliably flagging such failures without ground-truth labels remains
difficult. Token-level confidence conflates how a model phrases an answer with
whether it knows the answer, and uncertainty estimators that score a single
fixed response cannot observe how that answer would change under small
perturbations. We propose the \emph{Perturbation Reliability Score} (PRS), a
training-free uncertainty measure that probes reliability by perturbing a
question from two complementary directions---semantic rephrasing of the input
and low-rank Gaussian noise on the weights---and then freely regenerating a
small ensemble of responses. PRS reads three instability signals off this
ensemble: output fragmentation across final answers, answer drift toward
competing candidates, and local trace drift inside the reasoning process.
The score requires no labels, no fine-tuning, and no calibration on the
evaluation set. Across four mathematical reasoning benchmarks (Minerva,
MATH-500, GSM8K, DeepScaler) and two instruction-tuned models (Qwen2.5-3B and
Llama-3.2-1B), PRS consistently outperforms TokUR EU, a strong teacher-forced
epistemic-uncertainty baseline, improving answer-level detection AUROC by up to
$17.5$ points. A prefix re-query analysis further shows that PRS captures a
process-level failure pattern in which local token competition appears first,
provisional answers then shift, and visible fragmentation follows.
\end{abstract}


\section{Introduction}
\label{sec:intro}

Large language models are increasingly deployed as reasoning engines for
mathematics, science, and decision support, yet a fluent chain of thought is no
guarantee of a correct final answer. The same model that solves a competition
problem can, on a superficially similar problem, produce a confident derivation
that quietly drifts to the wrong result. For any system that acts on model
outputs, the practical question is therefore not only how often the model is
right, but whether we can tell---at test time and without access to the
ground-truth answer---which individual responses to trust. This is the problem
of answer-level reliability detection: assigning each response an uncertainty
score that ranks likely-wrong answers above likely-correct ones.

The dominant signal for this task has been the model's own probabilities, but
token-level confidence is a poor proxy for reasoning reliability. A model can be
locally fluent and assign high probability to every token of a wrong answer,
because softmax probabilities reflect how the model phrases its output far more
than whether the underlying conclusion is sound. Semantic-entropy approaches
address part of this by sampling multiple generations and clustering them by
meaning, so that uncertainty is measured over distinct answers rather than
surface forms. A second family treats uncertainty as epistemic and induces it
through the weights: perturbing a fixed response under Bayesian or low-rank
weight noise and measuring the resulting variance, as in token-level epistemic
estimators such as TokUR EU. Both directions improve over raw confidence, but
they share a structural limitation. They probe a single perspective---either the
temperature-sampling distribution of one model, or the weight-induced variance
along one frozen response---and they typically read uncertainty off a fixed
trajectory rather than observing how the answer itself reorganizes when the
problem is genuinely disturbed.

Our starting point is that an unreliable reasoning response should be
\emph{unstable}: if a model truly understands a problem, neither paraphrasing the
question nor jittering its weights should change the answer it commits to, nor
the candidates it entertains along the way. We operationalize this through the
\emph{Perturbation Reliability Score} (PRS), a training-free score that induces
instability from two independent sources and then lets the model regenerate
freely under each. Semantic rephrasings of the input test robustness to changes
in surface form that preserve the problem, while low-rank Gaussian weight
perturbations test robustness to small parameter changes; pooling the two axes
gives a compact ensemble of responses for the same question. Rather than scoring
a single fixed trace, PRS reads three complementary instability signals off this
ensemble---fragmentation of the final answers, drift of probability mass toward
competing answers, and local drift inside the reasoning trace---and combines
them into one score with fixed, label-free coefficients. The result is an
uncertainty estimator that observes how a response falls apart under
perturbation, instead of how confidently it was written.

We evaluate PRS as an answer-level reliability score on four mathematical
reasoning benchmarks (Minerva, MATH-500, GSM8K, and DeepScaler) and two
instruction-tuned models (Qwen2.5-3B and Llama-3.2-1B), comparing against TokUR
EU, a strong epistemic baseline that scores responses by teacher-forced
weight-perturbation variance. Across every shared setting PRS gives a stronger
reliability signal, improving AUROC by up to $17.5$ points and AUPRC by as much
as $36.5$ points, and the gains hold across both model families. A prefix
re-query analysis further suggests that PRS is not merely aggregating disagreement:
a subset of failures follows a consistent temporal pattern in which local
math-token competition emerges first, the model's provisional answer then shifts,
and visible response fragmentation appears only later. Our contributions are
threefold. First, we introduce a two-axis perturbation ensemble---semantic
rephrasing and low-rank weight noise---as a training-free probe of reasoning
reliability. Second, we define PRS, which combines output fragmentation, answer
drift, and trace drift into a single label-free score, and show that it
consistently beats a strong teacher-forced epistemic baseline across models and
datasets. Third, we provide a temporal analysis indicating that
perturbation-induced instability reflects a structured, process-level failure
mode rather than random noise.


\section{Method}
\label{sec:method}

We propose the \emph{Perturbation Reliability Score} (PRS), a training-free
uncertainty quantification method for reasoning responses. PRS estimates
reliability by inducing perturbation from two complementary sources, input
perturbations and weight perturbations, and then scoring the resulting responses
with three instability signals: output fragmentation, answer drift, and trace
drift.


\subsection{Perturbation Ensemble}
\label{sec:perturbation-ensemble}

PRS builds a perturbation ensemble from two independent perspectives, text
perturbations and weight perturbations. Text perturbations test whether the
model is robust to semantic changes of the input prompt, while weight
perturbations test whether the model is robust to small changes in its
parameters.

\paragraph{Text perturbation.}
Given a question $x$, we construct $R$ input variants
$\mathcal{X}(x)=\{\tilde{x}^{(r)}\}_{r=1}^{R}$.
Each $\tilde{x}^{(r)}$ is generated by an external rephrasing API under the
instruction to preserve the original problem's mathematical meaning,
constraints, and expected answer. This gives a text-side perturbation axis that
changes the surface form of the problem while keeping the target answer
unchanged.

\paragraph{Weight perturbation.}
In parallel, we construct $W$ model-side perturbations around the base model
$p_{\theta_0}$. We write
\[
\theta^{(w)} \sim \mathcal{Q}_{\rho}(\theta\mid \theta_0),
\qquad w=1,\ldots,W,
\]
where $\mathcal{Q}_{\rho}$ is centered at $\theta_0$ and $\rho$ controls the
perturbation strength. Following low-rank weight perturbation methods, we
perturb a fixed set of projection matrices $\mathcal{P}$ in the model. For a
selected matrix $W_\ell^0\in\mathbb{R}^{m_\ell\times n_\ell}$, let
$W_\ell^0 = U_\ell \operatorname{diag}(d_\ell)V_\ell^\top$
be its compact singular value decomposition. Let $U_{\ell,r'}$ denote the first
$r'$ columns of $U_\ell$, where $r'\ll r_\ell$. For the $w$-th perturbed model,
we sample
\[
E_{\ell}^{(w)}\in\mathbb{R}^{n_\ell\times r'},
\qquad
E_{\ell,ab}^{(w)}\sim\mathcal{N}(0,\rho_\ell^2),
\]
and define
$W_\ell^{(w)} = W_\ell^0 + U_{\ell,r'}\left(E_{\ell}^{(w)}\right)^\top$.
All parameters outside $\mathcal{P}$ remain unchanged, and the same perturbed
model instance $\theta^{(w)}$ is kept fixed throughout the generation of one
response.

\paragraph{Joint perturbation ensemble.}
We pool the two perturbation axes into a single ensemble rather than taking their
Cartesian product. The text axis contributes $R$ responses, each obtained by
decoding the base model $p_{\theta_0}$ on one rephrased prompt
$\tilde{x}^{(r)}$, and the weight axis contributes $W$ responses, each obtained
by decoding a perturbed model $p_{\theta^{(w)}}$ on the original prompt $x$. This
yields $K=R+W$ responses, which we index by $j\in[K]$ and write as
\[
\mathcal{O}(x)
=
\{o_j=(\tau_j,a_j,\mathbf{p}_j)\}_{j=1}^{K}.
\]
Here $\tau_j=(y_{j,1},\ldots,y_{j,T_j})$ is the reasoning trace, $a_j$ is the
extracted final answer, and $\mathbf{p}_j=\{p_{j,t}\}_{t=1}^{T_j}$ stores the
retained top-$M$ token distributions along the trace. Each response $j$ is
produced by a single model--prompt pair $(\theta_j,x_j)$, equal to either
$(\theta_0,\tilde{x}^{(r)})$ on the text axis or $(\theta^{(w)},x)$ on the weight
axis, so that
\[
p_{j,t}(v)
=
p_{\theta_j}(v\mid x_j,y_{j,<t}).
\]
Let $\mathcal{M}_{j,t}$ denote the retained top-$M$ token set at position $t$.
The PRS score below is computed entirely from this joint perturbation ensemble.



\subsection{Perturbation Reliability Score}
\label{sec:prs-score}

PRS combines three perturbation-induced signals,
\[
\mathrm{PRS}(x)
=
S_{\mathrm{out}}(x)
+
\lambda_a S_{\mathrm{ans}}(x)
+
\lambda_t S_{\mathrm{tr}}(x),
\]
where $S_{\mathrm{out}}$ measures output fragmentation, $S_{\mathrm{ans}}$
measures answer drift, and $S_{\mathrm{tr}}$ measures trace drift. The
coefficients $\lambda_a$ and $\lambda_t$ are fixed before evaluation and are not
learned from test labels.

\paragraph{Output fragmentation.}
Output fragmentation measures whether perturbations lead to inconsistent final
answers. Let $\phi(\cdot)$ be an answer canonicalization function that maps
equivalent answers to the same class; in practice $\phi$ applies answer
normalization, numeric tolerance, and symbolic simplification when applicable.
For response $j$, let $c_j=\phi(a_j)$ be its answer class, and let the empirical
distribution over observed answer classes be
\[
\widehat{p}_{\mathrm{cls}}(c\mid x)
=
\frac{1}{K}
\sum_{j=1}^{K}
\mathbf{1}\{c_j=c\}.
\]
We define output fragmentation as the probability mass that falls outside the
dominant answer class,
\[
S_{\mathrm{out}}(x)
=
1-\max_{c}\ \widehat{p}_{\mathrm{cls}}(c\mid x).
\]
This score is $0$ when all perturbation runs agree on a single answer class and
grows toward $1$ as the ensemble splits into multiple competing classes.

\paragraph{Answer drift.}
Output fragmentation only uses hard final answers. Answer drift instead measures
whether a response still assigns probability mass to competing answers while
producing its own final answer. Let $\mathcal{A}_j$ be the token positions of the
final-answer span in response $o_j$. For each observed answer class $c$, let
$\mathrm{Tok}(c)$ be the token ids obtained by tokenizing its canonical answer
string, and define the competing answer-token set as
$\mathcal{V}_{-j}=\bigcup_{c\neq c_j}\mathrm{Tok}(c)$.
The answer drift of response $j$ is
\[
b_j
=
\frac{1}{|\mathcal{A}_j|}
\sum_{t\in\mathcal{A}_j}
\sum_{v\in \mathcal{M}_{j,t}\cap \mathcal{V}_{-j}}
p_{j,t}(v),
\]
and the question-level answer drift is
$S_{\mathrm{ans}}(x)=K^{-1}\sum_{j=1}^{K}b_j$. A large $S_{\mathrm{ans}}$
indicates that the model commits to one answer while still placing
non-negligible probability on alternative candidates.

\paragraph{Trace drift.}
Trace drift measures local token-level instability inside the reasoning process.
Let $\mathcal{I}_j$ be the reasoning-token positions used for scoring, namely the
positions before the final-answer span. For each position $t\in\mathcal{I}_j$ we
define
\[
u_{j,t}
=
\sum_{v\in\mathcal{M}_{j,t}\setminus\{y_{j,t}\}}
p_{j,t}(v),
\]
the retained probability mass assigned to alternatives other than the generated
token. To avoid dilution by filler tokens and to keep the signal domain-agnostic,
we aggregate over the top-$\kappa$ fraction of positions with the highest local
spread, so that the trace drift of response $j$ is
\[
g_j
=
\frac{1}{|\mathcal{I}_j^{\kappa}|}
\sum_{t\in\mathcal{I}_j^{\kappa}}
u_{j,t},
\qquad
\mathcal{I}_j^{\kappa}\subseteq\mathcal{I}_j,
\]
where $\mathcal{I}_j^{\kappa}$ collects the highest-spread positions. The
question-level trace drift is $S_{\mathrm{tr}}(x)=K^{-1}\sum_{j=1}^{K}g_j$.

\paragraph{Final score.}
Substituting the three components above, PRS is
\[
\boxed{
\mathrm{PRS}(x)
=
\Big(1-\max_{c}\widehat{p}_{\mathrm{cls}}(c\mid x)\Big)
+
\lambda_a
\frac{1}{K}\sum_{j=1}^{K} b_j
+
\lambda_t
\frac{1}{K}\sum_{j=1}^{K} g_j
}.
\]
A larger PRS indicates lower reasoning reliability. The score is used directly
for uncertainty ranking and requires no supervised training or calibration on
the evaluation set.


\section{Experiments}
\label{sec:experiments}

We evaluate PRS as an answer-level reliability score for reasoning responses.
The goal is to determine whether an uncertainty score can distinguish incorrect
answers from correct ones without access to ground-truth labels at test time.
We first compare PRS with a strong epistemic-uncertainty baseline, and then
analyze whether the PRS signals reveal process-level failure patterns.

\subsection{Experimental Setup}
\label{sec:experimental-setup}

We evaluate on four mathematical reasoning benchmarks with different difficulty
profiles: Minerva, MATH-500, GSM8K, and DeepScaler. GSM8K consists primarily of
grade-school arithmetic problems, MATH-500 contains competition-style
mathematical reasoning questions, DeepScaler contains more challenging reasoning
problems, and Minerva provides a difficult benchmark on which models often
exhibit answer-level instability. For each question we use clean labels: a
response is treated as incorrect when its final answer is not equivalent to the
reference answer under a canonicalizing grader that normalizes fractions,
matrices, and LaTeX, and applies a relative numeric tolerance.

We conduct experiments with two open-source instruction-tuned models,
Qwen2.5-3B and Llama-3.2-1B. For every question PRS builds a pooled perturbation
ensemble of $K=8$ responses, drawn from $R=4$ text rephrasings of the prompt and
$W=4$ low-rank weight perturbations, all decoded greedily with a generation
budget of $2048$ tokens. The weight perturbations use rank $r'=4$ Gaussian noise
with strength $\rho=0.03$ on the query/key projection matrices, with the
perturbed model held fixed throughout each generation; the text rephrasings are
produced by an external rephrasing API. The PRS coefficients are fixed before
evaluation at $\lambda_a=0.05$ and $\lambda_t=0.03$ and are never tuned on test
labels, so the score is fully training-free.

We compare PRS against TokUR EU, a strong token-level uncertainty baseline based
on epistemic uncertainty under weight perturbations. TokUR EU scores a given
response through full-response teacher forcing, whereas PRS uses free perturbed
generations. This comparison highlights a practical distinction: TokUR EU
measures uncertainty along a fixed response, while PRS measures reliability from
the instability of responses induced by perturbations. We report AUROC and AUPRC
for answer-level reliability detection, treating incorrect answers as positive
examples. AUROC measures the overall ability of an uncertainty score to rank
incorrect answers above correct ones, while AUPRC emphasizes performance under
class imbalance; higher values indicate better reliability detection.

\subsection{Main Results}
\label{sec:main-results}

Table~\ref{tab:prs-vs-tokur-main} compares PRS with TokUR EU across the four
reasoning benchmarks, and PRS consistently outperforms TokUR EU on all shared
settings. On Qwen2.5-3B, PRS improves AUROC from $0.758$ to $0.822$ on Minerva,
from $0.801$ to $0.904$ on MATH-500, from $0.786$ to $0.961$ on GSM8K, and from
$0.712$ to $0.810$ on DeepScaler; the gains are even larger in AUPRC, where PRS
adds $19.7$ and $36.5$ points on MATH-500 and GSM8K respectively. The
improvements are not limited to Qwen2.5-3B: on Llama-3.2-1B, PRS also outperforms
TokUR EU on every dataset where TokUR EU is available, with AUROC gains of $12.2$
points on Minerva, $6.4$ points on MATH-500, and $10.3$ points on GSM8K, while
remaining strong on DeepScaler where TokUR EU does not produce a comparable
score. These results show that PRS provides a robust answer-level reliability
signal across model families and datasets, and they suggest that free
perturbation-induced response instability is a more direct indicator of answer
correctness than teacher-forced token-level epistemic uncertainty alone.

\begin{table*}[htbp]
  \centering
  \small
  \setlength{\tabcolsep}{5.0pt}
  \renewcommand{\arraystretch}{1.12}
  \caption{
    \textbf{Answer-level reliability detection: PRS vs. TokUR EU.}
    Metrics are AUROC / AUPRC under clean labels, with incorrect answers as the
    positive class. TokUR EU uses full-response teacher forcing, while PRS uses
    free perturbed generations. Best result in each column is in \textbf{bold}.
  }
  \label{tab:prs-vs-tokur-main}
  \resizebox{\textwidth}{!}{
  \begin{tabular}{l|cc|cc|cc|cc}
    \toprule
    \rowcolor{headerpink}
    & \multicolumn{2}{c|}{Minerva}
    & \multicolumn{2}{c|}{MATH-500}
    & \multicolumn{2}{c|}{GSM8K}
    & \multicolumn{2}{c}{DeepScaler} \\
    \rowcolor{headerpink}
    Method
    & AUROC & AUPRC
    & AUROC & AUPRC
    & AUROC & AUPRC
    & AUROC & AUPRC \\
    \midrule

    \multicolumn{9}{l}{\textbf{Qwen2.5-3B}} \\
    \rowcolor{rowgray}
    TokUR EU
      & 0.758 & 0.875
      & 0.801 & 0.659
      & 0.786 & 0.474
      & 0.712 & 0.782 \\

    PRS
      & \textbf{0.822} & \textbf{0.900}
      & \textbf{0.904} & \textbf{0.856}
      & \textbf{0.961} & \textbf{0.839}
      & \textbf{0.810} & \textbf{0.858} \\

    \midrule

    \multicolumn{9}{l}{\textbf{Llama-3.2-1B}} \\
    \rowcolor{rowgray}
    TokUR EU
      & 0.592 & 0.840
      & 0.727 & 0.842
      & 0.795 & 0.801
      & --- & --- \\

    PRS
      & \textbf{0.714} & \textbf{0.930}
      & \textbf{0.791} & \textbf{0.880}
      & \textbf{0.898} & \textbf{0.883}
      & 0.754 & 0.916 \\

    \bottomrule
  \end{tabular}
  }
\end{table*}

\subsection{Case Analysis}
\label{sec:temporal-mechanism}

Beyond static reliability ranking, we ask whether the PRS signals reveal a
process-level structure during reasoning. We therefore perform a prefix re-query
analysis and track three events along the normalized reasoning prefix: a
math-token spike $T_{\mathrm{math}}$, a prefix-answer distribution shift
$B_{\mathrm{shift}}$, and response fragmentation $A_{\mathrm{frag}}$. Here
$B_{\mathrm{shift}}$ is computed from adjacent prefix-answer distributions,
rather than from the entropy of a single prefix, so that it captures changes in
the model's provisional answer tendency as the reasoning trace unfolds.
Figure~\ref{fig:temporal-mechanism} summarizes the resulting temporal patterns.
Panel (a) shows that the strict event ordering $\tau_T < \tau_B < \tau_A$ is
enriched among wrong responses across datasets, and the enrichment is especially
strong on Minerva, where wrong responses exhibit a much higher frequency of this
ordering than both correct responses and the shuffle null. Panel (b) visualizes
the first event-time distributions on Minerva, where the three events occupy
different regions of the normalized prefix timeline, suggesting that the signals
are not merely simultaneous artifacts. Panel (c) further examines what happens
after the first prefix-answer shift: in wrong responses the probability of active
fragmentation increases after $\tau_B$, whereas correct responses remain
comparatively stable and eventually decrease.

These results suggest that PRS captures more than static response disagreement. A
subset of failures exhibits a process-level instability pattern in which local
math-token competition appears first, prefix-answer tendencies then shift, and
visible response fragmentation follows later. This supports the interpretation of
PRS as a layered reliability score that combines early trace-level instability
with later answer-level fragmentation.

\begin{figure*}[htbp]
  \centering
  \includegraphics[width=\textwidth]{mechanism.png}
  \caption{
    \textbf{Temporal mechanism analysis of perturbation-induced instability.}
    (a) Cross-dataset enrichment of the strict event ordering
    $\tau_T < \tau_B < \tau_A$ among wrong responses, correct responses, and a
    shuffle null.
    (b) First event-time distributions on Minerva for math-token spikes
    $T_{\mathrm{math}}$, prefix-answer shifts $B_{\mathrm{shift}}$, and response
    fragmentation $A_{\mathrm{frag}}$.
    (c) Post-shift fragmentation dynamics on Minerva: after the first
    prefix-answer shift $\tau_B$, wrong responses show increasing fragmentation
    activity, while correct responses remain lower and decrease over later
    stages.
  }
  \label{fig:temporal-mechanism}
\end{figure*}

\section{Conclusion}
\label{sec:conclusion}

We presented the Perturbation Reliability Score, a training-free uncertainty
measure that judges a reasoning response by how it destabilizes under two
complementary perturbations---semantic rephrasing of the input and low-rank
weight noise---rather than by the confidence with which it was written. By
reading output fragmentation, answer drift, and trace drift off a small pooled
ensemble, PRS turns response instability into a single label-free reliability
score. Across four mathematical benchmarks and two model families it
consistently outperforms a strong teacher-forced epistemic baseline, and a
temporal analysis indicates that the signal reflects a structured failure mode
rather than random noise. PRS currently targets math-style problems with
extractable final answers and relies on an external rephrasing step; extending it
to open-ended generation and to larger models, and replacing the rephrasing API
with an in-model perturbation, are natural directions for future work.


% \bibliography{iclr2026_conference}
% \bibliographystyle{iclr2026_conference}

% \appendix
% \section{Appendix}



\end{document}
